\documentclass[12pt]{article}
\usepackage[utf8]{inputenc}
\usepackage[russian]{babel}
\usepackage{amsmath, amssymb}
\usepackage{graphicx}
\usepackage{caption}
\usepackage{multicol}
\usepackage{array}
\usepackage{geometry}
\geometry{a4paper, margin=2cm}

\title{\vspace{-3em}\bfseries\LARGE Согревающие формулы}
\date{}

\begin{document}
	
	\maketitle
	
	\begin{multicols}{2}
		\noindent
		
		Зимой в студенческих общежитиях особенно остро ощущается нехватка горячего чая. Попробуем математически описать процесс получения "согревающего" напитка, используя некоторые простые формулы.
		
		\begin{figure}[H]
			\centering
			\includegraphics[width=0.4\linewidth]{tea1.jpg}
			\caption*{Рис. 1. Чашка с кипятком.}
		\end{figure}
		
		Предположим, что в одной чашке находится чай температурой $1^\circ C$, а в другой — $3^\circ C$. При смешении их поровну температура будет:
		\[
		t = \frac{1^\circ + 3^\circ}{2} = 2^\circ C.
		\]
		
		Добавим третью чашку также с температурой $1^\circ C$. Тогда средняя температура:
		\[
		t = \frac{1^\circ+3^\circ + 2\cdot1^\circ}{1+1} = \frac{5}{2}^\circ C.
		\]
		
		\begin{figure}[H]
			\centering
			\includegraphics[width=0.4\linewidth]{tea2.jpg}
			\caption*{Рис. 2. Смешивание жидкостей.}
		\end{figure}
		
		\noindent
		Можно обобщить процесс по следующей рекуррентной формуле:
		\[
		t_{k+1} = \frac{1}{n}(nt_k+1),
		\]
		где $n$ — количество порций, $t_k$ — текущая температура.
		
		\begin{center}
			\captionof{table}{Результаты подогрева при последовательных смешиваниях}
			\begin{tabular}{|c|c|c|c|}
				\hline
				\textbf{Шаг} & $t_k$ & $n$ & $t_{k+1}$ \\
				\hline
				1 & 1.0 & 1 & 2.0 \\
				2 & 2.0 & 2 & 2.5 \\
				3 & 2.5 & 3 & 2.75 \\
				4 & 2.75 & 4 & 2.875 \\
				\hline
			\end{tabular}
		\end{center}
		
		Таким образом, даже простое смешение жидкости можно описать при помощи несложной, но интересной математики.
		
	\end{multicols}
	
\end{document}
